\chapter{高阶微分与多元极值}
\label{chap:higher-differentials-extrema-solutions}

\begin{solution}{bilinear-operator-norm-bound}
设 $e_1,\ldots,e_n$ 和 $u_1,\ldots,u_m$ 分别是定义域两个因子中的标准基。
写成
\[
 h=\sum_{i=1}^nh_i e_i,\qquad k=\sum_{j=1}^mk_j u_j.
\]
由双线性性，
\[
 B(h,k)=\sum_{i,j}h_i k_jB(e_i,u_j).
\]
若 $M=\max_{i,j}\norm{B(e_i,u_j)}$，则 Cauchy--Schwarz 不等式给出
\[
 \norm{B(h,k)}
 \le M\Bigl(\sum_i|h_i|\Bigr)\Bigl(\sum_j|k_j|\Bigr)
 \le M\sqrt{nm}\,\norm h\norm k.
\]
这既给出了所需常数，也说明 $B$ 在原点连续。再由
\[
 B(h,k)-B(h_0,k_0)=B(h-h_0,k)+B(h_0,k-k_0)
\]
可知它在每一点连续。有限维性使有限组基向量足以控制全部取值；无限维中，
代数双线性映射不必自动连续。
\end{solution}

\begin{solution}{second-derivative-quadratic-map}
令 $q(x)=x^{\mathsf T}Ax$。直接展开
\[
 q(x+h)-q(x)=h^{\mathsf T}Ax+x^{\mathsf T}Ah+h^{\mathsf T}Ah,
\]
故
\[
 Dq(x)h=x^{\mathsf T}(A+A^{\mathsf T})h.
\]
把 $Dq(x)$ 对 $x$ 再微分，得到
\[
 D^2q(x)[h,k]=h^{\mathsf T}(A+A^{\mathsf T})k,
\]
它与 $x$ 无关。因此 $D^3q=0$，以后各阶微分也都为零。特别地，只有
$A$ 的对称部分 $(A+A^{\mathsf T})/2$ 对二次型有贡献。
\end{solution}

\begin{solution}{hessian-coordinate-representation}
若 $f$ 在 $a$ 二次可微，则对
$h=\sum_i h_i e_i$、$k=\sum_j k_j e_j$，双线性性给出
\[
 D^2f(a)[h,k]
 =\sum_{i,j}h_i k_jD^2f(a)[e_i,e_j]
 =h^{\mathsf T}H_f(a)k,
\]
其中 $(H_f(a))_{ij}=\partial_i(\partial_jf)(a)$：
第二个参数是 \(Df\) 所作用的向量，第一个参数是对 \(Df\) 求导的方向。
正文关于二阶 Fréchet 微分的对称性命题给出
$D^2f(a)[e_i,e_j]=D^2f(a)[e_j,e_i]$，因而 Hessian 对称。
这里不必要求 \(f\in C^2\)；\(C^2\) 是另一个便于检验的充分条件。
仅有全部二阶
偏导在一点存在，不足以保证两个混合偏导相等；选做题给出了反例。
\end{solution}

\begin{solution}{schwarz-rectangle-increment}
固定其余坐标，只写两个有关变量为 $s,t$。设
\[
 \Delta(h,k)=f(h,k)-f(h,0)-f(0,k)+f(0,0).
\]
对一元函数 $s\mapsto f(s,k)-f(s,0)$ 使用中值定理，存在
$\xi$ 位于 $0$ 与 $h$ 之间，使
\[
 \Delta(h,k)=h\{f_s(\xi,k)-f_s(\xi,0)\}.
\]
再对 $t\mapsto f_s(\xi,t)$ 使用中值定理，存在 $\eta$ 位于 $0$ 与 $k$
之间，使
\[
 \Delta(h,k)=hk (f_s)_t(\xi,\eta).
\]
交换两次作差的顺序，同样得到某个 $(\xi',\eta')$，使
\[
 \Delta(h,k)=hk (f_t)_s(\xi',\eta').
\]
当 $h,k\to0$ 且 $hk\ne0$ 时，四个中间点都趋于原点。若两个混合偏导在
邻域内存在且在原点连续，除以 $hk$ 后取极限即得
$(f_s)_t(0,0)=(f_t)_s(0,0)$。括号明确了先后求导的顺序。
平移坐标便得到任意点的结论。
\end{solution}

\begin{solution}{second-chain-rule}
令 $F=g\circ f$。一阶链式法则为
\[
 DF(x)=Dg(f(x))\,Df(x).
\]
沿方向 $k$ 对右端求导。第一个因子的变化由
$D^2g(f(x))[Df(x)k,\cdot]$ 给出，第二个因子的变化为
$D^2f(x)[k,\cdot]$。将所得线性算子作用于 $h$，便有
\[
 D^2F(x)[k,h]
 =D^2g(f(x))[Df(x)k,Df(x)h]
  +Dg(f(x))D^2f(x)[k,h].
\]
两个二阶微分对称时，右端也明显关于 $h,k$ 对称。
\end{solution}

\begin{solution}{radial-function-hessian}
记 $r^2=\norm{x}_2^2$。因为 $D(r^2)(x)h=2\langle x,h\rangle$，
链式法则给出
\[
 Df(x)h=2\phi'(r^2)\langle x,h\rangle,
 \qquad
 \nabla f(x)=2\phi'(r^2)x.
\]
再微分一次，得到
\[
 D^2f(x)[h,k]
 =4\phi''(r^2)\langle x,h\rangle\langle x,k\rangle
  +2\phi'(r^2)\langle h,k\rangle.
\]
因此
\[
 H_f(x)=4\phi''(r^2)xx^{\mathsf T}+2\phi'(r^2)I.
\]
\end{solution}

\begin{solution}{multi-index-count}
把 $k$ 个不可区分的单位分配给 $n$ 个位置，向量
$(\alpha_1,\ldots,\alpha_n)$ 对应在 $k$ 个单位中插入 $n-1$ 个隔板。
从 $n+k-1$ 个位置中选择 $n-1$ 个隔板位置，故共有
\(
 \binom{n+k-1}{n-1}=\binom{n+k-1}{k}
\)
个多重指标。二元三阶齐次项为
\[
 \frac16f_{xxx}h^3+\frac12f_{xxy}h^2k
 +\frac12f_{xyy}hk^2+\frac16f_{yyy}k^3,
\]
各偏导均在展开点取值。
\end{solution}

\begin{solution}{taylor-integral-remainder-derivation}
设 \(f\in C^r\)、\(r\ge1\)，令 $\gamma(t)=f(a+th)$。
记
\[
 I_j=\frac1{(j-1)!}\int_0^1(1-t)^{j-1}\gamma^{(j)}(t)\,dt
 \quad(1\le j\le r).
\]
一元微积分基本定理给出
\[
 \gamma(1)=\gamma(0)+\int_0^1\gamma'(t)\,dt.
\]
对 \(j\ge2\)，逐分量分部积分，得到
\[
 I_j=-\frac{\gamma^{(j-1)}(0)}{(j-1)!}+I_{j-1}.
\]
边界项在 \(t=1\) 为零，在 \(t=0\) 给出所列负项。
从 \(I_1=\gamma(1)-\gamma(0)\) 逐次代入，得到
\[
 \gamma(1)=\sum_{j=0}^{r-1}\frac{\gamma^{(j)}(0)}{j!}
 +\frac1{(r-1)!}\int_0^1(1-t)^{r-1}\gamma^{(r)}(t)\,dt.
\]
链式法则给出
$\gamma^{(j)}(t)=D^jf(a+th)[h^j]$，代回即为正文公式。
再分出常量 \(D^rf(a)\)，得到第 \(r\) 次项及 Peano 余项。
全过程只用到 \(r\) 阶连续导数与一元分部积分，不要求第 \(r+1\) 阶导数，
也不需要尚未建立的多重积分换序定理。
\end{solution}

\begin{solution}{taylor-multi-index-expansion}
对称 $k$ 线性映射 $D^kf(a)$ 作用在
$h=\sum_i h_i e_i$ 的 $k$ 个副本上。展开后，若指标 $i$ 出现
$\alpha_i$ 次，则相应项为
$\partial^\alpha f(a)h^\alpha$；这种排列共有 $k!/\alpha!$ 种。因此
\[
 D^kf(a)[h^k]
 =\sum_{|\alpha|=k}\frac{k!}{\alpha!}
   \partial^\alpha f(a)h^\alpha.
\]
两边除以 $k!$ 即得所需公式。
\end{solution}

\begin{solution}{taylor-exp-log-two-variables}
由
\[
 e^x=1+x+\frac{x^2}{2}+\frac{x^3}{6}+O(x^4),\qquad
 \log(1+y)=y-\frac{y^2}{2}+\frac{y^3}{3}+O(y^4),
\]
相乘并保留总次数不超过 $3$ 的项，得到
\[
 P_3(x,y)=y+xy-\frac12y^2+\frac12x^2y-\frac12xy^2+\frac13y^3.
\]
函数在原点附近属于 $C^4$，四阶 Taylor 余项为
$O((x^2+y^2)^2)$，它当然也是
$o((x^2+y^2)^{3/2})$。
\end{solution}

\begin{solution}{stationary-point-classification-polynomial}
驻点方程为
\[
 2x+y-2=0,\qquad x+2y+3=0,
\]
解得 $(x,y)=(7/3,-8/3)$。Hessian 为
\[
 H=\begin{pmatrix}2&1\\1&2\end{pmatrix},
\]
其特征值为 $1,3$，所以正定。该点是严格局部极小点。由于函数本身是
正定二次型加上线性项，它还是唯一的整体极小点。
\end{solution}

\begin{solution}{hessian-two-dimensional-criterion}
配方得到
\[
 ax^2+2bxy+cy^2
 =a\left(x+\frac ba y\right)^2+\frac{ac-b^2}{a}y^2
\]
（正定情形下先由取 $(1,0)$ 得 $a>0$）。故正定当且仅当
$a>0$ 且 $ac-b^2>0$。负定等价于 $-H$ 正定，即
$a<0$ 且 $ac-b^2>0$。若 $ac-b^2<0$，两个特征值异号，二次型不定。
行列式为零时存在非零零方向，不能由这个判据断定高阶项的符号。
\end{solution}

\begin{solution}{indefinite-hessian-saddle}
设 $Df(a)=0$，并取 $u,v$ 使
$D^2f(a)[u,u]>0$、$D^2f(a)[v,v]<0$。Taylor 展开为
\[
 f(a+tu)-f(a)=\frac{t^2}{2}D^2f(a)[u,u]+o(t^2),
\]
以及 $v$ 方向上的同类公式。第一式在充分小的非零 $t$ 下为正，第二式
为负。因而 $a$ 的每个邻域既含函数值较大的点，也含函数值较小的点，
$a$ 是鞍点。这一论证正是 Hessian 判别第三部分所依据的方向限制方法。
\end{solution}

\begin{solution}{compact-sphere-uniform-quadratic}
令 $S^{n-1}=\{u:\norm u=1\}$。连续函数 $u\mapsto Q(u)$ 在紧集
$S^{n-1}$ 上取得最小值。正定性说明这个最小值 $c$ 严格为正。于是对
$h\ne0$，
\[
 Q(h)=\norm h^2Q\left(\frac h{\norm h}\right)\ge c\norm h^2.
\]
有限维性用在单位球面的紧致性上。为看出它不能略去，取只有有限项
非零的实数列组成的空间，配以 \(\norm x_2=(\sum_jx_j^2)^{1/2}\)，并令
\[
 Q(x)=\sum_{j=1}^{\infty}\frac{x_j^2}{j}.
\]
对每个非零 \(x\)，\(Q(x)>0\)，且所有和均为有限和。
对应双线性型由 Cauchy--Schwarz 不等式满足
\(|\sum_jx_jy_j/j|\le\norm x_2\norm y_2\)，所以连续。
但标准单位向量 \(e_j\) 满足 \(Q(e_j)=1/j\to0\) 且
\(\norm{e_j}_2=1\)，不可能有统一正下界。
\end{solution}

\begin{solution}{local-minimum-hessian-semidefinite}
若 $a$ 是局部极小点，对任意 $h$ 考察
$\varphi(t)=f(a+th)$。它在 $0$ 取得局部极小，故
$\varphi''(0)=D^2f(a)[h,h]\ge0$。所以 Hessian 半正定。
反向命题不成立。例如 $f(x,y)=x^2-y^4$ 在原点的 Hessian 为
$\operatorname{diag}(2,0)$，半正定；但沿 $x=0$ 有 $f(0,y)=-y^4<0$，
沿 $y=0$ 有 $f(x,0)=x^2>0$，原点是鞍点。
\end{solution}

\begin{solution}{quartic-extrema-trichotomy}
$x^4+y^4\ge0$，且等号只在原点成立，故原点是严格（事实上整体）极小点。
对其相反数，原点是严格整体极大点。对 $x^4-y^4$，沿 $y=0$ 取正值，
沿 $x=0$ 取负值，所以原点是鞍点。三者在原点的一、二、三阶微分都相同，
差别首次出现在四阶齐次项。
\end{solution}

\begin{solution}{same-hessian-different-behavior}
两个函数在原点的梯度都为零，Hessian 都是
$\operatorname{diag}(2,0)$。然而
$x^2+y^4>0$ 对每个非零 $(x,y)$ 成立，原点是严格极小点；
$x^2-y^4$ 沿两坐标轴取相反符号，原点是鞍点。Hessian 的核是 $y$ 轴，
四阶项在这个核上的符号决定了两者的不同。
\end{solution}

\begin{solution}{first-nonzero-homogeneous-term}
设
$f(a+h)-f(a)=P_k(h)+o(\norm h^k)$，且 $P_k$ 在非零点为正。
齐次性表明 $P_k(-h)=(-1)^kP_k(h)$，所以 $k$ 不可能为奇数。
在单位球面上，$P_k$ 连续且严格为正，故有最小值 $c>0$。于是
\[
 P_k(h)\ge c\norm h^k.
\]
取 $h$ 充分小，使余项绝对值不超过 $c\norm h^k/2$，便有
$f(a+h)>f(a)$。若 $P_k$ 在单位球面上分别于 $u,v$ 取正、负值，则沿
$tu,tv$ 的余项分别被固定符号的主项控制，任意邻域内都有函数值在
$f(a)$ 两侧的点，故 $a$ 是鞍点。
\end{solution}

\begin{solution}{odd-leading-term-saddle}
非零齐次多项式 $P_k$ 至少在某个单位向量 $u$ 上非零。若 $k$ 为奇数，
$P_k(-u)=-P_k(u)$，因此它在单位球面上取正负值。将上一题的第二部分
用于 Taylor 展开的首个非零齐次项，便知相应驻点是鞍点。
\end{solution}

\begin{solution}{flat-function-no-finite-taylor-test}
记 $r=\sqrt{x^2+y^2}$。在穿孔平面上，任意阶偏导均为
\[
 P(x,y,r^{-1})e^{-1/r^2}
\]
的形式，其中 \(P\) 为多项式：求导时
\(\partial_x r^{-q}=-qx r^{-q-2}\)，
\(\partial_x e^{-1/r^2}=2xr^{-4}e^{-1/r^2}\)，对 \(y\) 同理，
所以该形式归纳保持。它在坐标轴的非原点处也有定义。
在 \(0<r\le1\) 上，对每个多重指标 \(\alpha\)，相应偏导 \(g_\alpha\)
可估为 \(C_\alpha r^{-N_\alpha}e^{-1/r^2}\)。
指数函数快于任意幂的估计说明，对每个 \(L>0\)，
\(g_\alpha(x,y)=o(r^L)\)。

将每个 \(g_\alpha\) 在原点赋值零。它们不仅连续延拓，而且确实仍是
延拓函数的偏导：归纳假定 \(\partial^\alpha f=g_\alpha\)，则在原点
沿坐标方向 \(e_i\) 的差商满足
\[
 \left|\frac{g_\alpha(te_i)-g_\alpha(0)}t\right|\longrightarrow0.
\]
因此该偏导等于零，即 \(g_{\alpha+e_i}(0)\)；非原点处则由普通求导成立。
从 \(\alpha=0\) 开始归纳，所有阶偏导均存在且连续，故 \(f\in C^\infty\)，
每阶 Taylor 多项式都为零。另一方面，$r>0$ 时
$e^{-1/r^2}>0=f(0,0)$，所以原点仍是严格局部极小点。这个例子说明
有限阶乃至全部 Taylor 系数都不必决定一般 $C^\infty$ 函数的局部值。
\end{solution}

\begin{solution}{mixed-partials-fail-without-continuity}
定义
\[
 f(x,y)=
 \begin{cases}
 \displaystyle\frac{xy(x^2-y^2)}{x^2+y^2},&(x,y)\ne(0,0),\\[4pt]
 0,&(x,y)=(0,0).
 \end{cases}
\]
公式在两坐标轴上自动给出零。固定 $y$，由定义计算
\[
 f_x(0,y)=\lim_{h\to0}\frac{f(h,y)-f(0,y)}h
 =\lim_{h\to0}y\frac{h^2-y^2}{h^2+y^2}=-y.
\]
故 $(f_x)_y(0,0)=-1$。同理
\[
 f_y(x,0)=\lim_{k\to0}x\frac{x^2-k^2}{x^2+k^2}=x,
\]
所以 $(f_y)_x(0,0)=1$。两个混合偏导都存在但不相等；Schwarz 定理中的
“在一点连续”的条件不能只用“在一点存在”代替。由正文的二阶
Fréchet 对称性命题，这个函数在原点也不二次 Fréchet 可微。
\end{solution}
